% =====================================================================
%  sec6.tex — 6. Four-Dimensional Gauge Theory And Stacks（含 6.1 Stacks）
%  原文 pp.10–14
% =====================================================================
\section{四维规范场论与叠}

对物理学家来说，在研究涉及规范场论的对偶时，从四维开始是自然的，因为四维常被发现是规范场论对偶的自然舞台。理由很简单：曲率——规范场论中最基本的概念之一——是一个 2-形式；在 $d$ 维中，2-形式的对偶是 $(d-2)$-形式，所以只有当 $d=4$ 时它仍然是 2-形式\zhnt{6.1}。这提示 $d=4$ 是「一个规范场论的对偶可能仍是规范场论」的最自然的维数。

此外，$\mathcal{N}=4$ 超对称 Yang--Mills 理论（最初由 Brink 等 (1977) 构造）是自然的起点，因为它具有可能的最大超对称，并拥有可追溯到 Montonen 与 Olive (1977) 早期工作的著名对偶\zhnt{6.2}。结果确实如此：几何 Langlands 在四维 $\mathcal{N}=4$ 超 Yang--Mills 理论的一个扭转版本中有其自然的起源。这种扭转与导致 Donaldson 理论的 $\mathcal{N}=2$ 超 Yang--Mills 理论的扭转相当类似。

那个特定的动机对某些人来说可能晦涩；因此，为理解几何 Langlands，我在这里改用另一种方式来解释为什么从四维开始是自然的，而不是只依赖 $\MH(G,C)$ 的 B-模型与 A-模型。

首先，任何空间 $X$ 的 B-模型与 A-模型，都是支配映射 $\Phi:\Sigma\to X$ 的量子 $\sigma$ 模型的扭转版本，其中 $\Sigma$ 是一个二维流形（更确切地说，是一个玻色维数为 2 的超流形）\zhnt{6.3}。由于 A-模型在其最基本的形式中涉及对满足适当条件的全纯映射 $\Phi:\Sigma\to X$ 的计数，$\Sigma$ 与 $\Phi$ 在 A-模型中的角色是清楚的。镜像对称表明，用映射 $\Phi:\Sigma\to X$ 来表述 B-模型也必定是正确的，而物理学家在通常的表述中正是这样做的。

在目前的情形中，我们粗略地说对某个紧 Lie 群 $G$ 与二维流形 $C$ 的空间 $\MH(G,C)$ 的 B-模型与 A-模型感兴趣。因此粗略地说，我们想研究映射 $\Phi:\Sigma\to\MH(G,C)$ 的 $\sigma$ 模型，其中如前 $\Sigma$ 是一个辅助的二维流形。

这一描述之所以粗略，是因为 $\MH(G,C)$ 有奇点\yz{此外，这些奇点比轨形奇点更坏。轨形奇点不会造成任何困难。关于几何 Langlands 中轨形奇点的讨论，见 Frenkel 与 Witten (2007)。}，因此以 $\MH(G,C)$ 为目标的 $\sigma$ 模型并非真正良定义。所以，一个完整的描述不能纯粹用「目标空间是视为抽象流形的 $\MH(G,C)$」的 $\sigma$ 模型来给出。我们需要一个更完整的描述，告诉我们如何处理这些奇点。这会是什么呢？

按定义，$\MH(G,C)$ 中的一个点在规范等价下决定一个满足 Hitchin 方程的对 $(A,\phi)$。$A$ 与 $\phi$ 是定义在 $C$ 上的场，所以我们把它们更明确地写成 $(A(y),\phi(y))$，其中 $y$ 是 $C$ 上的坐标。

现在设有一个映射 $\Phi:\Sigma\to\MH(G,C)$，其中 $\Sigma$ 是带局部坐标 $x$ 的黎曼面。这样的映射由一对同时依赖 $x$ 的 $(A(y),\phi(y))$ 描述。于是我们可以用同时依赖 $x$ 与 $y$ 的场 $(A(x,y),\phi(x,y))$ 来描述映射 $\Phi$。我们希望把这些场解释为四维流形 $M=\Sigma\times C$ 上的场\zhnt{6.4}。对 $(A(x,y),\phi(x,y))$ 并不完全是一组 $M$ 上的自然场，但可以自然地补全为这样一组。例如，$A(x,y)$ 局部地是切于 $M=\Sigma\times C$ 中第二个因子的 1-形式；要得到四维规范场，我们应放松 $A$ 切于第二个因子这一条件。类似地，我们可以把 $\phi$ 延拓为 $\Sigma\times C$ 上取值于伴随丛的 1-形式。$\mathcal{N}=4$ 超 Yang--Mills 理论——更确切地说，与几何 Langlands 相关的它的扭转版本——正是通过以最小方式把这组场补全为一个超对称组合而得到的。

在 $\mathcal{N}=4$ 超 Yang--Mills 理论中，没有类似于 $\MH(G,C)$ 的奇点。例如，规范场的空间是仿射空间，其他场（如 $\phi$）取值于线性空间。使「映射 $\Phi:\Sigma\to\MH(G,C)$ 的 $\sigma$ 模型」难以定义的那些奇点问题，在 $M=\Sigma\times C$（或任何其他四维流形）上定义规范场论时都没有类似物。两者的关系是：二维 $\sigma$ 模型是四维规范场论的近似\zhnt{6.5}。当遇到 $\MH(G,C)$ 的奇点时，这一近似失效。任何涉及这些奇点的问题，都应在底层的四维规范场论中处理。

不过，远离奇点时，只需考虑描述映射 $\Phi:\Sigma\to\MH(G,C)$ 的那组较小的场就够了。许多问题不依赖奇点，对这些问題，用二维 $\sigma$ 模型与镜像对称来描述是足够的。

\subsection{叠}

作为结论，我们将与上述讨论在通常的数学理论中的对应物建立联系。我们从丛而不是 Higgs 丛开始，因为这一情形更容易解释。

在通常的数学理论中，几何 Langlands 对应的右手边是用 D-模来描述的，粗略地说，是在黎曼面 $C$ 上所有全纯 $G_{\C}$-丛的模空间上。然而，人们考虑的不是半稳定全纯 $G_{\C}$-丛 $E\to C$ 的模空间 $M(G,C)$，而是所有这类丛的「叠」$\Bun_{G}(C)$ 上的 D-模\zhnt{6.6}。主要原因是：要定义 Hecke 算子的作用，就必须允许不稳定丛。不稳定丛与 $M(G,C)$ 的非轨形奇点有关。

什么是叠？粗略地说，它是一个可以处处局部地描述为商的空间。平凡的情形是：一个叠实际上可以整体地描述为商。把 $\Bun_{G}(C)$ 解释为整体商，意味着要找一对 $(Y,W_{\C})$，由一个光滑代数簇 $Y$ 与一个作用在 $Y$ 上的复 Lie 群 $W_{\C}$ 组成，具有以下性质：全纯 $G_{\C}$-丛 $E\to C$ 的同构类应与 $Y$ 上的 $W_{\C}$-轨道一一对应；并且对每个 $E\to C$，其自同构群应同构于 $W_{\C}$ 中固定 $Y$ 中对应点的子群。

如果要求 $Y$ 与 $W_{\C}$ 是有限维的，那么以此方式表示叠 $\Bun_{G}(C)$ 的对 $(Y,W_{\C})$ 并不存在\zhnt{6.7}。确实，$G_{\C}$-丛 $E\to C$ 可以任意不稳定，所以其自同构群的维数没有上界。因此没有有限维的 $W_{\C}$ 能包含所有这些自同构群作为子群。

然而正如 Atiyah 与 Bott (1982) 所证明的，为简单起见取 $G$ 为伴随型，存在自然的无穷维对 $(Y,W_{\C})$。只需取 $Y$ 为给定（最初只在拓扑上定义的）$G$-丛 $E\to C$ 上所有联络的空间。定义 $W$ 为丛 $E$ 的所有规范变换构成的群；于是若 $E$ 拓扑平凡，可把 $W$ 等同于 $\Maps(C,G)$。然后取 $W_{\C}$ 为 $W$ 的复化，换言之 $\Maps(C,G_{\C})$。（这个复化如下作用在 $Y$ 上：把联络 $A$ 对应到相应的 $\bar\partial$ 算子 $\bar\partial_{A}$；于是复值规范变换按 $\bar\partial_{A}\to g\bar\partial_{A}g^{-1}$ 作用。）

设我们被摆在这样的问题面前：已知 $M(G,C)$ 的奇点，要弄清映射 $\Phi:\Sigma\to M(G,C)$ 的超对称 $\sigma$ 模型的含义。（这是我们真正问题——涉及 $\MH(G,C)$ 而不是 $M(G,C)$——的一个练习情形。）代数几何方面的朋友会告诉我们用叠 $\Bun_{G}(C)$ 替换 $M(G,C)$。我们把这一叠解释为对 $(Y,W_{\C})$，其中 $Y$ 是 $G$-丛 $E\to C$ 上所有联络的空间，$W_{\C}$ 是复化的规范变换群。叠的连通分支对应 $E$ 的拓扑选择。

所谓以一对 $(Y,W_{\C})$（$W_{\C}$ 是作用在复流形 $Y$ 上的复 Lie 群）为目标的超对称 $\sigma$ 模型，在有限维情形下，我们指的是规范群为 $W$（$W_{\C}$ 的极大紧子群）、目标为 $Y$ 的规范不变超对称 $\sigma$ 模型。实际上，要定义这个 $\sigma$ 模型，我们希望 $Y$ 是一个带 $W$-不变（当然不是 $W_{\C}$-不变）Kähler 结构的 Kähler 流形。$\sigma$ 模型作用量含有一项，它是 $W$ 作用的矩映射的平方。这一项恰在矩映射为零时取极小。把矩映射置零并对 $W$ 取商，这二者合起来在经典上等价于对 $W_{\C}$ 取商\zhnt{6.8}。

要写下作用量中涉及矩映射平方的那一项（事实上，要写下规范场的动能项），需要在 $W$ 的李代数上有一个不变正定二次型。若 $W$ 是有限维的，这种形式的存在等价于 $W$ 是紧的\zhnt{6.9}。然而在 $W=\Maps(C,G)$（对某个空间 $C$）的无穷维情形中，相应的二次型也存在。（$W$ 的李代数的一个元素是 $C$ 上取值于 $\mathfrak{g}$ 的函数 $\epsilon$，二次型由 $\int_{C}d\mu\,(\epsilon,\epsilon)$ 定义，其中 $(\ ,\ )$ 是 $\mathfrak{g}$ 上不变正定二次型，$\mu$ 是 $C$ 上适当的测度。）

现在，设我们构造从黎曼面 $\Sigma$ 到叠 $\Bun_{G}(C)$ 的映射的二维 $\sigma$ 模型，叠按上述方式理解。这个 $\sigma$ 模型的规范变换群是什么？一般地，在任何空间 $\Sigma$ 上规范群为 $W$ 的规范场论中，规范变换群（为简单起见，取拓扑平凡 $W$-丛的）是从 $\Sigma$ 到 $W$ 的映射群，即 $\Maps(\Sigma,W)$。在我们的情形，$W$ 又是 $\Maps(C,G)$。所以 $\Maps(\Sigma,W)$ 就是 $\Maps(M,G)$，其中 $M=\Sigma\times C$。但这正是 $M$ 上规范群为 $G$ 的规范场论中的规范变换群。在目前情形，$\Sigma$ 与 $C$ 是二维流形，$M$ 是四维流形。我们已经到达四维规范场论。若把定义再追问得细一些，我们就会知道：映射 $\Phi:\Sigma\to\Bun_{G}(C)$ 的超对称 $\sigma$ 模型应被理解为四维流形 $M=\Sigma\times C$ 上规范群为 $G$ 的四维 $\mathcal{N}=2$ 超对称规范场论。（这是扭转后与 Donaldson 理论相关的那套理论。）

现在回到原来的问题。几何 Langlands 对偶是关于 $\MH(G,C)$（相应的 Higgs 丛模空间）及其对偶群版本 $\MH(\LG,C)$ 的 B-模型与 A-模型的陈述，而不是关于 $M(G,C)$ 的。为处理奇点，我们想把这一情形「叠化」。我们现在处在超凯勒的语境中，叠的恰当概念应把这一点纳入。（代数几何家所谓的 Higgs 丛的叠并不足以胜任这一情形，因为它过于强调一个复结构。）由于四元数 Lie 群不存在，我们不能要求把 $\MH(G,C)$ 构造为光滑空间被四元数 Lie 群取商\zhnt{6.10}。然而，辛商的概念在超凯勒世界中有很好的类比，即 Hitchin 等 (1987) 所描述的超凯勒商。对我们为 $M(G,C)$ 所解释的东西的类比是：把 $\MH(G,C)$ 实现为光滑空间 $Y$ 被群 $W$ 的超凯勒商。用有限维的 $Y$ 与 $W$ 来做这件事可能做不到，但在无穷维世界中，这个问题在 Hitchin 关于 Hitchin 方程的原始论文 Hitchin (1987a) 中有一个自然的解。（$Y$ 是 $C$ 上对 $(A,\phi)$ 的空间，$W=\Maps(C,G)$。）以此为输入，并解释「目标为对应于 $\MH(G,C)$ 的超凯勒叠的 $\sigma$ 模型」应当意味着什么，人们就得到 $\mathcal{N}=4$ 超 Yang--Mills 理论的扭转版本——而这也正是 Kapustin 与 Witten (2007) 的出发点。
